\chapter*{Abstract}
\addcontentsline{toc}{chapter}{Abstract}

\vspace{0.5cm}

The exponential growth of connected devices in 5G and beyond networks demands innovative multiple access techniques to maximize spectral efficiency while ensuring fairness. Power-Domain Non-Orthogonal Multiple Access (PD-NOMA) has emerged as a promising solution, allowing multiple users to share identical time-frequency resources through power-domain multiplexing and Successive Interference Cancellation (SIC). However, the critical challenge of intelligent user pairing—determining which users should share resources—significantly impacts system performance and becomes computationally prohibitive for large-scale deployments.

This project presents a comprehensive framework combining traditional optimization-based clustering methods with modern deep learning techniques for NOMA user pairing. We implement and evaluate four classical strategies: Static Pairing, Balanced Pairing, Blossom Maximum-Weight Matching, and Bipartite Graph Matching, using realistic Channel State Information (CSI) generated from the standardized 3GPP TR 38.901 Urban Macro (UMa) channel model.

Building upon these baselines, we propose \textbf{NOMA-GNN}, a novel Graph Neural Network architecture based on GraphSAGE that treats user pairing as a graph representation learning problem. The framework employs multi-task learning to jointly predict: (1) pairing feasibility with 92.5\% AUC, (2) achievable sum-rates with MAE of 0.031 bits/Hz, and (3) optimal power allocation coefficients with MAE of 0.009. Physical constraints including SIC feasibility (8 dB threshold) and angular separation (25°) are enforced during candidate edge generation.

Extensive experimental evaluation on 50,000 synthetic scenarios demonstrates that NOMA-GNN achieves 95.8\% of optimal Blossom throughput while providing a 20× computational speedup (6.4 ms vs 127.6 ms for N=12 users) with O(N log N) complexity compared to O(N³) optimization methods. The model exhibits robust generalization to unseen user densities, channel conditions, and cell geometries with minimal performance degradation (<4\%). Furthermore, NOMA-GNN achieves the highest fairness among all methods (Jain's index: 0.903) and demonstrates practical deployability with a lightweight 520 KB model suitable for edge base stations.

The complete implementation includes: (1) realistic 3GPP-compliant channel generation incorporating path loss, shadowing, and fading; (2) graph-based pairing algorithms with computational complexity analysis; (3) end-to-end deep learning pipeline with data preprocessing, training, and real-time inference; (4) comprehensive visualization and evaluation framework. Results indicate that data-driven GNN approaches can effectively replace computationally expensive optimization while maintaining near-optimal performance, enabling practical real-time NOMA deployment in dense urban environments.

\vspace{0.5cm}

\noindent\textbf{Keywords:} Non-Orthogonal Multiple Access (NOMA), Graph Neural Networks (GNN), GraphSAGE, User Pairing, Power Allocation, Multi-Task Learning, 3GPP Channel Model, Successive Interference Cancellation (SIC), Deep Learning for Wireless Networks, 5G/6G Systems.

\newpage
